\chapter{Álgebras de Boole Finitas}

En el capítulo anterior descubrimos que la estructura algebraica pura de Boole encierra en su interior una rigurosa topología de orden. Comprobamos cómo todo par de elementos está sometido a las leyes del supremo ($\vee$) y el ínfimo ($\wedge$) dentro de un retículo acotado, distributivo y complementado.

Esta dualidad álgebra/topología es la clave matemática que nos va a revelar la forma exacta de todas las álgebras de Boole que tienen un número finito de elementos.

\section{Estructura Atómica}

Para comprender la anatomía de un álgebra finita, debemos identificar sus "ladrillos fundamentales", es decir, los elementos indivisibles a partir de los cuales se puede construir todo el conjunto operando con supremos (sumas).

\begin{quote}\textbf{Definicion (Átomos):}\label{atomos}
Un elemento $x \in \mathbb{B}$ se denomina \textbf{átomo} si es un elemento estrictamente positivo, $x \ne \bot,$ y no existe ningún elemento intermedio entre él y el $\bot.$ Formalmente:
\[ \mathit{atom}(x) \iff (x > \bot) \wedge \left( \forall y \in \mathbb{B}, \bot < y \le x \implies y = x \right) \]
\end{quote}

En términos puramente algebraicos, un átomo $x$ es aquel cuyo producto (ínfimo) con cualquier otro elemento $y \in \mathbb{B}$ es el aniquilador total o bien él mismo (absorbente total):
\[ \mathit{atom}(x) \iff (x \ne \bot) \wedge \left( \forall y \in \mathbb{B}, (x \wedge y = \bot) \vee (x \wedge y = x) \right) \]
El conjunto de todos los átomos de un álgebra de Boole se denota como $\mathit{Atom}(\mathbb{B}).$
Una propiedad inmediata es que el ínfimo de dos átomos distintos es siempre nulo:
\[ \forall a, b \in \mathit{Atom}(\mathbb{B}), a \ne b \implies a \wedge b = \bot \]

\begin{quote}\textbf{Definicion (Hiperátomos (Co-átomos)):}\label{hiperatomos}
De forma dual, un \textbf{hiperátomo} o \textbf{co-átomo} es un elemento estrictamente inferior a $\top$ tal que no existe ningún elemento intermedio entre él y el máximo.
\[ \mathit{hatom}(x) \iff (x < \top) \wedge \left( \forall y \in \mathbb{B}, x \le y < \top \implies y = x \right) \]
\end{quote}

En cualquier álgebra finita no trivial (donde $\top \ne \bot$), los conjuntos $\mathit{Atom}(\mathbb{B})$ e $\mathit{Hatom}(\mathbb{B})$ nunca están vacíos.

\section{El Teorema de Representación de Stone (Caso Finito)}

Si tomamos el conjunto de todos los átomos $\mathit{Atom}(\mathbb{B}),$ podemos generar el conjunto potencia $\wp(\mathit{Atom}(\mathbb{B})),$ es decir, el conjunto de todos sus posibles subconjuntos. Vamos a construir una función $\varphi$ que conecte este álgebra de subconjuntos con el álgebra de Boole original.

\subsection{El Funtor $\varphi$}
Definimos la función $\varphi : \wp(\mathit{Atom}(\mathbb{B})) \to \mathbb{B}$ que toma un subconjunto de átomos $S$ y devuelve el supremo topológico (la suma booleana) de todos ellos:
\[
\varphi(S) = \begin{cases}
\bot & \text{si } S = \emptyset \\
\bigvee_{a \in S} a & \text{si } S \ne \emptyset
\end{cases}
\]

Esta función preserva de forma natural las operaciones del álgebra de subconjuntos hacia las del álgebra de Boole original:
\begin{align*}
\varphi(S_1 \cup S_2) &= \varphi(S_1) \vee \varphi(S_2) \\
\varphi(S_1 \cap S_2) &= \varphi(S_1) \wedge \varphi(S_2) \\
\varphi(\mathit{Atom}(\mathbb{B}) \setminus S) &= \neg \varphi(S)
\end{align*}

\subsection{Isomorfismo y Cardinalidad}
El gran triunfo matemático para las álgebras finitas se resume en el siguiente teorema.

\begin{quote}\textbf{Teorema (Representación de Álgebras de Boole Finitas):}\label{rep_stone}
Toda álgebra de Boole finita $\mathbb{B}$ es algebraicamente isomorfa al álgebra del conjunto potencia de sus átomos. Es decir, la función $\varphi$ es una biyección perfecta:
\[ \mathbb{B} \simeq \wp(\mathit{Atom}(\mathbb{B})) \]
\end{quote}
\begin{proof}
Dado que el álgebra es finita, no existen cadenas infinitas descendentes. Todo elemento $x \in \mathbb{B}$ (salvo el $\bot$) acota por arriba a al menos un átomo. Por distributividad y ortogonalidad de los átomos ($a \wedge b = \bot$), cualquier elemento $x$ puede expresarse de forma única como el supremo de los átomos que lo preceden: 
\[ x = \bigvee \{ a \in \mathit{Atom}(\mathbb{B}) \mid a \le x \} \]
Por lo tanto, la función inversa de $\varphi$ está perfectamente definida, demostrando que $\varphi$ es sobreyectiva e inyectiva (biyectiva).
\end{proof}

De este isomorfismo se desprende una consecuencia colosal, formulada previamente de soslayo en nuestro análisis de los conjuntos topológicos.

\begin{quote}\textbf{Teorema (Cardinalidad de un Álgebra Finita):}\label{cardinalidad}
Si un álgebra de Boole es finita, su número total de elementos debe ser, forzosamente, una potencia de 2.
\[ |\mathbb{B}| = 2^n \quad \text{donde} \quad n = |\mathit{Atom}(\mathbb{B})| \]
\end{quote}

Cualquier conjunto que no tenga exactamente $2, 4, 8, 16, \ldots$ elementos \textbf{jamás podrá} constituir un álgebra de Boole, sin importar qué operaciones intentemos definir sobre él. Y aún más importante: todas las álgebras de Boole finitas que tengan el mismo número de elementos son exactamente la misma álgebra (son algebraicamente isomorfas). Solo hay \textit{un} álgebra de Boole de 2 elementos, \textit{una} de 4 elementos, \textit{una} de 8 elementos, etc.

\section{Generalización a Álgebras Infinitas}

En el caso de las álgebras de Boole infinitas, el Teorema de Representación que acabamos de ver no se cumple de forma incondicional. Como comprobamos en el primer capítulo con los modelos matemáticos, en el infinito la topología puede volverse mucho más exótica. 

Para que un álgebra infinita sea isomorfa al conjunto potencia de sus átomos, debe poseer dos propiedades estructurales estrictas:
\begin{enumerate}
    \item \textbf{Ser atómica:} Todo elemento no nulo de $\mathbb{B}$ debe estar acotado inferiormente por al menos un átomo.
    \item \textbf{Ser completa:} Todo subconjunto (incluso infinito) de elementos debe poseer un supremo y un ínfimo que pertenezcan obligatoriamente al álgebra.
\end{enumerate}

\begin{quote}\textbf{Teorema (Isomorfismo de Álgebras Atómicas y Completas):}\label{rep_stone_infinito}
Cualquier álgebra de Boole $\mathbb{B}$ que sea simultáneamente completa y atómica es isomorfa al álgebra del conjunto potencia de sus átomos:
\[ \mathbb{B} \simeq \wp(\mathit{Atom}(\mathbb{B})) \]
\end{quote}

Si falla alguna de estas dos propiedades, el isomorfismo se rompe. Podemos revisitar los modelos topológicos de cardinalidad $\aleph_0$ para ilustrar este fenómeno:
\begin{itemize}
    \item El \textbf{Ejemplo 9} (la familia de subintervalos racionales) es un álgebra infinita que carece por completo de átomos. Puesto que cualquier subintervalo racional siempre puede subdividirse en dos más pequeños, no existe el "ladrillo indivisible".
    \item El \textbf{Ejemplo 10} (los subconjuntos finitos y cofinitos de $\mathbb{N}$) \textit{sí} es un álgebra atómica (cuyos átomos son los conjuntos unitarios $\{n\}$). Sin embargo, \textit{no} es completa. Si reuniésemos una infinidad de estos átomos (por ejemplo, los números pares), su supremo lógico sería el conjunto de todos los pares, pero dicho conjunto no es finito ni cofinito y, por ende, no pertenece al álgebra. Al no ser completa, esta álgebra no es isomorfa al conjunto potencia $\wp(\mathbb{N}).$
\end{itemize}

Esta diversidad matemática infinita es fascinante, pero en la electrónica digital nos da luz verde para restringir nuestros esfuerzos puramente al caso finito y bivaluado. Sabemos de antemano que, bajo limitaciones finitas, cualquier álgebra de Boole es idéntica al espacio vectorial $n$-dimensional $\mathbb{B}_2^n,$ un tema que exploraremos a fondo en el siguiente capítulo.
